\documentclass[11pt]{article}
\usepackage{booktabs}
\usepackage[margin=1in]{geometry}
\usepackage{fontspec}
\setmainfont{lmroman10-regular.otf}[BoldFont=lmroman10-bold.otf,ItalicFont=lmroman10-italic.otf,BoldItalicFont=lmroman10-bolditalic.otf]
\setsansfont{lmsans10-regular.otf}[BoldFont=lmsans10-bold.otf,ItalicFont=lmsans10-oblique.otf]
\setmonofont{lmmono10-regular.otf}[ItalicFont=lmmono10-italic.otf]
\title{Xe Lm Fontspec}
\author{A. Author}
\date{1 January 2026}
\begin{document}
\maketitle
\begin{abstract}
This synthetic benchmark document exercises the engine with typical content.
\end{abstract}
\tableofcontents
\section{Large Variance}\label{sec:1}

The structural impact varies the clear stability of the comparison. In the structural comparison, the equilibrium shapes a possible stability and the market. In the marginal signal, the market varies a possible state and the portfolio. We find that the sufficient pattern determines the price, while the early population conditions the early return. We find that the primary pattern drives the system, while the broad experiment shapes the early feature. A average context changes each policy in the general source.

We find that the complex framework varies the process, while the optimal risk predicts the observed design. In the stable feature, the noise improves a partial noise and the comparison. We find that the stable evidence shapes the estimate, while the primary pattern changes the broad profit. We find that the primary size changes the system, while the marginal return supports the primary response. A strong behavior varies each gain in the partial control.

\section{Sufficient Income}\label{sec:2}

A narrow judgment improves each strategy in the constant variance. The typical balance tracks the optimal system of the finding. We find that the complex finding drives the judgment, while the local approach relates the complex utility. The dynamic framework bounds the robust range of the value. A marginal feature reflects each decision in the strong judgment. A robust forecast drives each approach in the average loss (see Section~\ref{sec:7}).

We find that the strong sector supports the result, while the structural strategy relates the local decision. The efficient design improves the joint index of the capital. We find that the common balance changes the hypothesis, while the joint regime conditions the large signal. The simple pattern conditions the implicit stability of the experiment (see Section~\ref{sec:7}). In the general constraint, the pressure shapes a broad loss and the limit.

\section{Empirical Pattern}\label{sec:3}

In the possible impact, the error relates a complex hypothesis and the investment (see Section~\ref{sec:6}). A average dynamic determines each threshold in the observed labor. The implicit scale reduces the primary volume of the cost. We find that the possible loss predicts the decision, while the implicit context reflects the efficient condition. The empirical parameter stabilizes the primary observation of the scale. A simple shock reduces each layer in the narrow process (see Section~\ref{sec:7}).

\begin{table}[tbp]\centering\caption{High policy summary.}\label{tab:b3}\begin{tabular}{lrrr}\toprule Coefficient & Factor & Utility & Coefficient \\ \midrule
coefficient & 5.19 & 16.69 & 38.43 \\
risk & 30.57 & 85.12 & 7.71 \\
function & 36.97 & 73.09 & 32.51 \\
investment & 53.11 & 30.27 & 14.68 \\
selection & 38.59 & 42.35 & 83.14 \\
outcome & 46.41 & 44.13 & 13.66 \\
demand & 80.40 & 74.00 & 81.87 \\
weight & 93.27 & 7.16 & 69.55 \\
\bottomrule\end{tabular}\end{table}

Table~\ref{tab:b3} lists the values. We find that the observed decision affects the pattern, while the narrow function changes the sufficient stability. The persistent state reduces the central yield of the outcome (see Section~\ref{sec:1}).

In the joint supply, the constraint predicts a central curve and the hypothesis. The final strategy varies the typical pressure of the dynamic. A partial approach drives each data in the direct impact. A empirical portfolio bounds each distribution in the local design (see Section~\ref{sec:1}). The central policy determines the observed curve of the factor.

\section{Primary Regime}\label{sec:4}

We find that the simple dynamic relates the efficiency, while the complex profit tracks the observed cluster. We find that the complex policy relates the flow, while the clear delay reduces the common portfolio (see Section~\ref{sec:2}). A narrow context shapes each threshold in the global decision. We find that the efficient dynamic tracks the framework, while the early yield determines the robust population. A modest constraint improves each finding in the partial balance. The efficient dynamic affects the partial source of the interval.

The benchmark edit marker reads BENCH-A.

The typical income reduces the large strategy of the labor. The random return predicts the modest investment of the pressure. A efficient process reflects each production in the efficient rate. A implicit pressure tracks each information in the implicit condition. The large uncertainty bounds the large choice of the process.

\section{Simple Value}\label{sec:5}

We find that the sharp efficiency affects the control, while the complex risk tracks the partial variance. In the partial yield, the policy reduces a global interaction and the interval. We find that the persistent knowledge conditions the production, while the early interval tracks the empirical utility. We find that the global feature reflects the margin, while the optimal regime tracks the robust experiment (see Section~\ref{sec:2}). The implicit effect predicts the final production of the decision (see Section~\ref{sec:5}). We find that the persistent control explains the data, while the persistent income reduces the global outcome.

We find that the sufficient forecast reflects the stability, while the joint system reflects the broad judgment. In the typical trend, the measure supports a large result and the measure (see Section~\ref{sec:6}). We find that the structural noise reduces the decision, while the possible policy bounds the optimal latency. In the simple noise, the approach explains a optimal pattern and the trade. A large gain changes each test in the persistent supply.

\section{Stable Production}\label{sec:6}

A stable trade predicts each interaction in the structural function. A general noise shapes each production in the partial margin. We find that the strong layer predicts the trade, while the clear capital stabilizes the persistent measure. In the general judgment, the effect improves a modest curve and the system. The final signal drives the sharp return of the feature. The observed balance limits the modest growth of the method.

\begin{table}[tbp]\centering\caption{Early scale summary.}\label{tab:b6}\begin{tabular}{lrrr}\toprule Error & Labor & Choice & Margin \\ \midrule
signal & 26.37 & 59.52 & 74.19 \\
demand & 30.84 & 76.79 & 96.95 \\
example & 94.17 & 0.88 & 50.23 \\
constraint & 43.52 & 26.55 & 33.15 \\
flow & 64.94 & 87.92 & 12.00 \\
pressure & 53.51 & 18.72 & 45.43 \\
margin & 67.48 & 75.31 & 41.38 \\
selection & 80.47 & 63.33 & 71.10 \\
\bottomrule\end{tabular}\end{table}

Table~\ref{tab:b6} lists the values. In the central capital, the experiment reduces a dynamic context and the result (see Section~\ref{sec:5}). We find that the direct measure varies the information, while the robust equilibrium predicts the high probability.

A efficient control reflects each cluster in the possible pressure. In the low labor, the rate explains a empirical interval and the hypothesis (see Section~\ref{sec:2}). We find that the robust curve shapes the labor, while the low constraint conditions the simple choice. In the random context, the approach reflects a global cost and the uncertainty. A strong regime varies each variable in the modest prediction (see Section~\ref{sec:8}).

\section{Random Equilibrium}\label{sec:7}

The sharp approach relates the sharp volume of the margin. A high equilibrium conditions each stability in the possible policy. The direct assumption supports the observed production of the regime. We find that the primary test reduces the measure, while the random information stabilizes the high size. The efficient gain reflects the stable weight of the state. We find that the clear judgment relates the regime, while the modest experiment improves the early yield.

The sufficient finding predicts the central capital of the latency. We find that the marginal regression predicts the index, while the typical method supports the final efficiency (see Section~\ref{sec:7}). We find that the sharp balance reflects the price, while the implicit probability supports the sharp dynamic. A complex curve shapes each condition in the constant technique. The strong limit reflects the random trend of the analysis.

\section{Persistent Judgment}\label{sec:8}

The large strategy reduces the clear framework of the analysis. The global hypothesis relates the active spread of the selection. The high labor relates the observed method of the probability. The primary signal improves the possible estimate of the strategy. The broad measure reflects the narrow experiment of the control. A direct knowledge drives each design in the broad pressure.

A low evidence bounds each measure in the central capital. We find that the marginal behavior bounds the regression, while the persistent analysis conditions the early network. We find that the direct latency reflects the trend, while the structural interval conditions the constant uncertainty. In the strong evidence, the interaction improves a efficient error and the sector. A empirical loss shapes each variable in the local scale (see Section~\ref{sec:5}).

\end{document}
